\documentclass[12pt]{article}
\usepackage{graphicx} % Required for inserting images
\usepackage{fancyhdr}
\usepackage{hyperref}

\pagestyle{fancy}
\fancyhf{}
\lhead{Final Assignment}

\title{Computer Workshop Final Assignment}
\author{Habibollah Panbechi 401411345}

\begin{document}

\begin{titlepage}
    \maketitle
\end{titlepage}
\tableofcontents
\newpage

\section{Git and GitHub}

\subsection{Repository Initialization and Commits}
\begin{enumerate}
    \item First I created a new repository by clicking the new button on the GitHub homepage.
    \item Then created a new local repository with git on my computer by entering ``git init'' in terminal.
    \item Created a basic ``READEME.md'' and pushed it on GitHub.
    \item Cloned ``\href{https://github.com/MiliAxe/CW1402Final}{https://github.com/MiliAxe/CW1402Final}'' to use it's workflow. copied ``.github'' folder and all of it's components in my project directory. Then added it to git by using this command ``git add .github'' and pushed to GitHub again.
    \item Created a basic \LaTeX{} file called \emph{main.tex}, added to git like before the pushed again. This time I created a new tag ``v.0.0.0'' and assigned it to the lastest commit by using this command ``git tag -a v0.0.0 \verb|<|HASHOFCOMMIT\verb|>|''.
    \item Now the repository is set. Every time I finish a subsection, I commit and push it with a new tag, then GitHub automatically complies this document.
\end{enumerate}


\subsection{GitHub Actions for LaTeX Compilation}
To use a GitHub workflow you have to setup a .yml file and configure it your desires.
\begin{enumerate}
    \item Create a directory named ``.github'', by using ``mkdir .github'' command.
    \item Create another directory inside that named ``workflows''.
    \item Create a file name ``main.yml'' inside that by using ``touch main.yml''.
    \item Configure your .yml file to your desires.
\end{enumerate}
My configuration compiles my \LaTeX{} document whenever I push something with the tag *.*.* (Where * is any sequence of characters).

\newpage

\section{Exploration Tasks}
\subsection{Discovering Vim's Advanced Features}
\begin{enumerate}
    \item \textbf{Marks:} Marks in Vim allow you to quickly jump to specific locations within a file. the `m' key is used to mark the cursor's position, followed by a letter to name it. for example `ma' marks the current spot. then by pressing `a' you can easily jump to it. a mark can be deleted by ``:delmarks a'' or ``:delmarks!'' which deletes all marks.
    \item \textbf{Folding:} In every IDE you can fold a piece of code to focus your attention on something specific. This feature is available in vim. To manually create a fold, you can use visual mode to select text and then type `zf'. `zo' opens a fold, `zc' closes one, `zM' closes all folds and 'zR' opens all of them.
    \item \textbf{Persistent Undo and Redo:} By using ``:undofile'' you can browse your undo history and navigate through it. Also by adding ``set undofile'' to your .vimrc file your can undo and redo even after closing your file. Because it makes them persistent. 
\end{enumerate}

\subsection{Memory profiling}
\subsubsection{Memory Leak}
Memory leaks happen when you dynamically allocate memory and do not free it when you are done using it. (As simple as that!) The problem is after your function is finished executing, you do not have any access to that chunk and therefore cannot free it. This can cause reduced performance of you program. tools like \emph{Valgrind} can help detecting and debugging this kind of issue. 

\subsubsection{Memory profilers}
Valgrind is a programming tool for memory debugging, memory leak detection, and profiling. Valgrind actively monitors memory allocation and deallocation during the program's execution. If it detects memory leaks (blocks of memory that are not freed), it reports the details, including the file and line number where the memory was allocated. It can also be used for general program profiling, function call frequencies, and execution time of different code sections.

\subsection{GNU/Linux Bash Scripting}
\subsubsection{fzf}
\begin{itemize}
    \item Fuzzy searching is a technique used in text matching, where the goal is to find approximate matches for a given search query rather than requiring an exact match. When fuzzy searching a string S, it compares to all other strings and gives each one a degree of similarity. Two popular string metrics are Levenshtein distance and Jaro-Winkler distance. fzf is a handy CLI tool to fuzzy search on your Unix machine.
    \item This command uses `ls' to list files and directories. Then pipes the result to `fzf', so the user can interactively search the result.
\end{itemize}

\subsubsection{Using fzf to find your favorite PDF}
\begin{enumerate}
    \item I prefer to use the find command, because I think it's more general and fd is only a user-friendlier version of it. Anyway the command to find all the directories of files with pdf extension is ``find / -type f -iname "*.PDF"''.
    \item Use this command to interactively select your pdf and open it, with your default pdf viewer``interactively''.
\end{enumerate}

\subsubsection{Opening the file using Zathura}
Zathura's handy syntax let us execute the following command: \newline``zathura \$(find / -type f -iname "*.PDF")''.

\newpage

\section{Git and FOSS}
\subsection{Issues}
\begin{figure}
    \centering
    \includegraphics{Screenshot 2024-01-27 225143.png}
    \caption{screenshot of my issue}
\end{figure}

\end{document}

